\chapter{LASA 层自适应脉动阵列架构}
\label{chap:architecture}

本文的目标不是单独实现一个卷积算子，而是在边缘 MPSoC 上完成从 DDR 图像输入、PS 侧调度、PL 侧卷积计算、OFM 回写到 YOLO decode 的完整链式推理。为此，本文提出一种层自适应脉动阵列加速器（Layer-Adaptive Systolic Accelerator, LASA）。LASA 以统一 INT8 systolic array 为计算基础，通过 layer descriptor 和 AXI-Lite 配置支持不同卷积核、输出通道分块、K pass、量化参数、激活和池化组合，从而在同一套硬件上覆盖 YOLOv3-tiny Conv0--Conv9 的多种层形状。

\section{系统设计目标与软硬件分工}

真实 YOLOv3-tiny 链式推理对加速器提出三类要求。第一，PL 侧卷积核心需要保持稳定的 INT8/INT32 计算语义，使 Conv0--Conv9 的输出能够与 golden 逐层比较。第二，PS/PL 边界需要承载真实层的数据搬运和背压，而不是依赖小规模仿真中的理想输入。第三，系统必须暴露足够的计数器和调试接口，使端到端延迟可以被分解为 DMA、IFM 重排、阵列计算、partial-PSUM 反馈、OFM 后处理和软件 decode 等阶段。

因此，本文采用“PS 负责任务组织，PL 负责规则密集计算”的协同方式。PS 侧 A53 程序维护 layer descriptor、DDR 缓冲区、DMA 启停、输出重排和 YOLO decode；PL 侧 LASA 接收寄存器配置和 AXI-Stream 数据，在片上完成卷积、重量化、激活、池化和 OFM packet 输出。第~\ref{chap:dataflow} 章讨论的 KCS 执行模型以及 BSD、OCRR、OPF-P 三类优化策略均建立在这一软硬件边界之上。

\section{MPSoC 总体结构}

图~\ref{fig:system_overview} 给出系统级结构。DDR 中保存输入图像包、量化权重、中间特征图和输出 packet；PS A53 通过 AXI-Lite 配置 LASA，并通过四路 AXI DMA 搬运 bias、weight、IFM 和 OFM 数据。LASA 以 AXI-Stream 接收三路输入流，以 AXI-Stream 输出 OFM debug packet 到 DDR，随后由软件重排并执行 YOLO decode。

\begin{figure}[htb]
\centering
\includegraphics[width=\textwidth]{\detokenize{img/picture/MPSOC架构3.png}}
\caption{MPSoC 端到端系统结构}
\label{fig:system_overview}
\end{figure}

该结构分离控制面和数据面。AXI-Lite 只承载层参数、tile 参数、量化参数、stream 配置和性能计数器；大规模张量数据全部通过 AXI DMA 搬运。这一划分使 PL 内部数据流可以持续演进，而外部软件 ABI 和 DMA 拓扑保持稳定。

\section{LASA 数据通路概览}

图~\ref{fig:accelerator_datapath} 展示 LASA 的数据通路。输入侧包含 bias loader、weight loader 和 IFM feeder；中间是 layer scheduler、systolic array 和 partial-PSUM 存储；输出侧包含 collector/drain、requant、activation、pooling 和 OFM writer。该图只描述各功能块的连接关系，具体的 KCS 分块、片上重排、PSUM 反馈和 pass 级预取在第~\ref{chap:dataflow} 章展开。

\begin{figure}[htb]
\centering
\includegraphics[width=\textwidth]{\detokenize{img/picture/LASA数据通路.png}}
\caption{LASA 数据通路概览}
\label{fig:accelerator_datapath}
\end{figure}

LASA 的最终板级配置采用 ROWS=18、COLS=8、$Cout_{tile}=16$。阵列沿 K 维接收 18 路 IFM lane；每个列位置包含两个 INT8 output-channel lane，因此 COLS=8 对应 16 个并行输出通道。两个 lane 的乘积分别累加到独立 INT32 partial sum 中，并在 final pass 后进入逐通道重量化路径。该 dual-lane INT8 column 设计提高了输出通道维度并行度，也决定了 PSUM packet、requant lane 和 OFM writer 的基本粒度。

\section{层自适应机制}

LASA 的“层自适应”体现在四个方面。第一，计算路径原生支持 3x3 与 1x1 卷积。对 3x3 卷积，IFM feeder 根据 $oy,ox,k_y,k_x,c_i$ 生成展开后的 K 维向量；对 1x1 卷积，硬件使用 native vector path，避免将 1x1 强行映射为稀疏 3x3 后产生无效 K pass。这样 Conv7 和 Conv9 等检测头可以使用更短的 $K_{total}$。

第二，后处理路径采用可配置旁路。final K pass 输出首先进入逐通道重量化模块；重量化后的 uint8 输出进入 activation LUT，可根据层描述选择 identity、ReLU/Leaky ReLU 等语义。若当前层启用 pooling，则 OFM pooling 模块执行 $2\times2$、stride=2 的 maxpool；若不启用 pooling，则数据直接旁路到 OFM byte stream。该机制使同一套硬件可以覆盖 YOLOv3-tiny 中带 pool 和不带 pool 的层。

第三，层参数由 descriptor 驱动。软件通过 layer descriptor 配置输入/输出尺寸、kernel、stride、padding、$K_{total}$、Cout blocks、spatial tile、量化参数、activation 和 pool 标志。硬件不需要为不同层重新综合，而是通过 AXI-Lite 寄存器和 stream packet 边界切换执行语义。

第四，数据流模式可组合。LASA 的基础数据通路可以与 batch stream、raw-HWC cache、partial-PSUM overlap、continuous collector 和 during-compute prefetch 等模式组合使用。对于以 1x1/3x3 标准卷积、逐层激活和池化为主的轻量化边缘 CNN，该机制具备扩展到其他顺序型卷积层的基础；对 depthwise/group convolution、residual add、concat 等结构，则需要后续扩展专门的数据路径或软件融合策略。

\section{链式推理执行流程}

一次 Conv0--Conv9 推理从 DDR 中的图像包开始。软件读取或准备输入图像，按 layer descriptor 写入当前层的尺寸、量化、tile 和 stream 配置，随后启动 OFM S2MM DMA 以及 bias、weight、IFM 三路 MM2S DMA。LASA 在当前层结束后输出带地址的 OFM packet，软件将其重排为 HWC buffer，并作为下一层输入或最终 YOLO decode 输入。

当前 OFM 输出采用 debug packet 语义：每个有效输出 byte 携带地址并写入 DDR。该格式不是最高带宽效率的连续 HWC burst，但保留了地址信息，便于板级阶段进行逐层重排、golden 比对和错误定位。第~\ref{chap:runtime} 章将进一步说明 runtime、寄存器接口、验证模式和性能计数器如何围绕该路径组织。

\section{本章小结}

本章提出并定义了 LASA 层自适应脉动阵列架构。LASA 通过 dual-lane INT8 systolic array、1x1/3x3 原生计算路径、可旁路后处理和 descriptor-driven 配置，为 YOLOv3-tiny 链式推理提供统一硬件执行基础。下一章将在 LASA 上进一步提出 KCS 三维分块执行模型，并围绕 PS/PL 交互、IFM 数据排布和 K-pass 边界气泡展开优化。
